%!TEX encoding = UTF8
%!TEX root = 

%%%%%%%%%%%%%%%%%%%%%%%%
%%%%%%%%%% SET 1 %%%%%%%%%%
%%%%%%%%%%%%%%%%%%%%%%%% 



\exe{origin=set1,}{
	Soient $f(x) = 3+2x$ et $g(x) = x^2 + 2x - 1$ deux polynômes à coefficients réels.
	Calculer
	\begin{multicols}{4}
	\begin{enumerate}
		\item $f(x)+g(x)$ ;
		\item $f(x) g(x)$ ;
		\item $f(x)^2$ ; et
		\item $g(f(x))$.
	\end{enumerate}
	\end{multicols}
}{exe:sum-product}{
	\begin{enumerate}
		\item 
			\[  f(x)+g(x) = 3+2x + x^2 + 2x - 1 = x^2 + 4x + 2. \]
		\item 
			\[ f(x) g(x) = (3+2x)(x^2 + 2x - 1) = 3x^2 + x - 3. \]
		\item 
			\[ f(x)^2 = (3+2x)^2 = 4x^2 + 12x + 9. \] 
		\item 
			\[ g(f(x)) = f(x)^2 + 2f(x) - 1 = (4x^2 + 12x + 9) + 2(3+2x) - 1 = 4x^2 + 16x + 14. \]
	\end{enumerate}
}

\exe{origin=set1,}{
	Pour chaque polynôme de $\R[x]$ suivant, donner : son degré ; son coefficient dominant ; et son coefficient constant. Dire si le polynôme est unitaire ou non.
	\begin{multicols}{2}
	\begin{enumerate}[itemsep=5pt]
		\item
		$p_1(x) = x^3 + x - 1$
		\item
		$p_2(x) = -x^7 + 6x + 4$
		\item
		$p_3(x) = x^3 + x^5 - x$
		\item
		$p_4(x) = 1 + x^2 - x^4 - x^3 + x$
	\end{enumerate}
	\end{multicols}
	\emph{Notation : le degré d'un polynôme $p$ est noté $\deg(p)$.}
}{exe:degre-unitaire}{
	\begin{enumerate}
		\item
		$\deg(p_1) = 3$. Le coefficient dominant est $1$ donc $p_1$ est unitaire.
		Le coefficient constant est $-1$.
		\item
		$\deg(p_2) = 7$. Le coefficient dominant est $-1$ donc $p_2$ n'est pas unitaire.
		Le coefficient constant est $4$.
		\item
		$\deg(p_3) = 5$. Le coefficient dominant est $1$ donc $p_3$ est unitaire.
		Le coefficient constant est $0$.
		\item
		$\deg(p_4) = 4$. Le coefficient dominant est $-1$ donc $p_4$ n'est pas unitaire.
		Le coefficient constant est $1$.
	\end{enumerate}
}


\exe{origin=set1,difficulty=1}{
	Calculer de tête le coefficient
	\begin{multicols}{2}
	\begin{enumerate}[itemsep=5pt]
		\item
		en $x$ de $(x+3)(x-4)$ ;
		\item
		en $x^2$ de $(1+x+x^2+x^3)(2x^2 + 3x - 5)$ ; \columnbreak
		\item
		en $x^6$ de $\left(3-4x-x^3+7x^4 + 2x^5\right)^2$ ;  et
		\item
		en $x^{50}$ de \mbox{$(1 + x^{25} + x^{50} + x^{100})(1 - x^{25} - x^{50} - x^{100})$}.
	\end{enumerate}
	\end{multicols}
}{exe:coeff-reading}{
	\begin{enumerate}
		\item
		Seuls les produits $x\times(-4)$ et $3\times x$ contribuent au terme en $x$.
		Le coefficient en $x$ est donc $-4+3 = -1$.
		
		On pourra vérifier que $(x+3)(x-4) = x^2 -x - 12$.
		\item
		Seuls les produits $1\times 2x^2$ ; $x \times 3x$ ; et $x^2 \times(-5)$ contribuent au terme en $x^2$.
		Le coefficient en $x^2$ est donc $2+3-5 = 0$.
		
		On pourra vérifier que $(1+x+x^2+x^3)(2x^2 + 3x - 5) = 2x^5 + 5x^4 - 2x - 5$.
		\item
		Remarquons que $\left(3-4x-x^3+7x^4 + 2x^5\right)^2 = \left(3-4x-x^3+7x^4 + 2x^5\right)\left(3-4x-x^3+7x^4 + 2x^5\right)$.
		Ainsi, seuls les produits $-4x\times2x^5$ ; $-x^3 \times (-x^3)$ ; et $2x^5 \times(-4x)$ contribuent au terme en $x^6$.
		Le coefficient en $x^6$ est donc $-8 + 1 - 8 = -15$.
		
		On pourra vérifier que $\left(3-4x-x^3+7x^4 + 2x^5\right)^2 = 4x^{10} + 28x^9 + 45x^8 -14x^7 - 15x^6 - 44x^5 -6x^3 + 16x^2 -24x + 9$ (merci Wolfram !).
		\item
		Seuls les produits $1\times(-x^{50})$ ; $x^{25} \times (-x^{25})$ ; et $x^{50} \times 1$ contribuent au terme en $x^{50}$
		Le coefficient de $x^{50}$ est donc $-1 -1 +1 = -1$.
		
		On pourra vérifier que $(1 + x^{25} + x^{50} + x^{100})(1 - x^{25} - x^{50} - x^{100}) = -x^{200} - 2x^{150} -2x^{125} -x^{100} - 2x^{75} -x^{50} + 1$.
	\end{enumerate}
}

\exe{origin=set1,difficulty=1}{
	Calculer le coefficient 
	%\begin{multicols}{2}
	\begin{enumerate}[itemsep=5pt]
		\item
		en $x^6$ de $(3x^2 + 7x - 3)(7x - 2x^4 + 3)$
		\item
		constant de $\left(17 + 31320x - 10^{31}x^2 - \frac13 x^7\right)\left(41x^2 + 76x - 2\right)$ ;
		\item
		en $x$ de $(x+3)(x-4)(x-8)$ ;
		\item
		en $x^2$ de $(x+1)(x-1)(x+3)(x+6)(x-2)$ ;
		\item
		en $x^6$ de $(1+x)^7$.
	\end{enumerate}
	%\end{multicols}
	\emph{Il est possible de vérifier ses résultats à l'aide de WolframAlpha (\href{https://www.wolframalpha.com/input?i=\%28x\%2B1\%29\%28x-1\%29\%28x\%2B3\%29\%28x\%2B6\%29\%28x-2\%29}{https://www.wolframalpha.com/}).}
}{exe:coeff-reading-3}{
	\begin{enumerate}
		\item
		Seul le produit $3x^2 \times (-2x^6)$ contribue au terme en $x^6$.
		Le coefficient en $x^6$ est donc $-6$.
		\item
		Seul le produit $17\times(-2)$ contribue au terme constant.
		Le coefficient constant est donc $-34$.
		\item
		Trois produits contribuent au terme en $x$ : parmi les trois parenthèses, il faut choisir $x$ une fois, et la constante les deux autres fois.
		Les produits contribuant au terme en $x$ sont donc $x\times(-4)\times(-8)$ ; $3\times x\times(-8)$ ; et $3\times(-4)\times x$.
		Le coefficient en $x$ est donc $32-24-12 = -4$.
		
		On pourra vérifier que $(x+3)(x-4)(x-8) = x^3 -9x^2-4x+96$.
		\item
		en $x^2$ de $(x+1)(x-1)(x+3)(x+6)(x-2)$ ;
		Il s'agit ici de choisir $x$ deux fois parmis les cinq parenthèses.
		Il existe ${5 \choose 2} = 10$ façons de faire. Celles-ci sont
		\begin{multicols}{2}
		\begin{enumerate}[label=Produit \arabic* :]
			\item
			$x\times x \times 3\times6\times(-2) = -36x^2$ ;
			\item
			$x\times(-1)\times x \times6\times(-2) = 12x^2$ ;
			\item
			$x\times(-1)\times3\times x\times(-2) = 6x^2$ ;
			\item
			$x\times(-1)\times3\times6\times x = -18x^2$ ;
			\item
			$1\times x\times x \times 6\times(-2) = -12x^2$ ;
			\item
			$1\times x \times 3 \times x \times(-2) = -6x^2$ ;
			\item
			$1\times x \times 3 \times6\times x = 18x^2$ ;
			\item
			$1\times(-1)\times x \times x \times (-2) = 2x^2$ ;
			\item
			$1\times(-1)\times x \times6\times x = -6x^2$ ; et
			\item
			$1\times(-1)\times3\times x \times x = -3x^2$. (ouf !)
		\end{enumerate}
		\end{multicols}
		Le coefficient en $x^2$ est donc $-36+12+6-18-12-6+18+2-6-3 = -43$. 
		
		Notons que le correcteur s'y est pris à deux fois pour obtenir le bon résultat, vérifié par WolframAlpha qui fournit $(x+1)(x-1)(x+3)(x+6)(x-2) = x^5 +7x^4-x^3 -43x^2 +36$.
		\item
		Comme $(1+x)^7 = (1+x)(1+x)(1+x)(1+x)(1+x)(1+x)(1+x)$, pour obtenir $x^6$, il faut choisir $x$ six foix parmi les sept parenthèses, et 1 une fois.
		Il existe 7 façons de faire, le choix de $x$ pouvant être fait en première, deuxième, troisième, ..., ou septième parenthèse.
		Le coefficient en $x^6$ est donc 7.
		
		On pourra vérifier douloureusement que $(1+x)^7 = x^7 +7x^6 + 21x^5 + 35x^4 + 35x^3 + 21x^2 + 7x + 1$.
		Question ouverte : pourquoi les coefficients forment-ils un palindrome ?
	\end{enumerate}
}

\exe{origin=set1,}{
	Montrer que le coefficient constant d'un polynôme $p$ est égal à $p(0)$.
	
	Quel est le coefficient constant de $p(x) = (x-1)^{231} - (x-1)^{436}$ ?
}{exe:coeff-const}{
	Si $p(x) = a_0 + a_1x + a_2x^2 + \dots + a_d x^d$ où $d=\deg(p)$, alors $p(0) = a_0 + 0 + 0 + \dots + 0 = a_0$, le coefficient constant.
	
	Le coefficient constant de $p(x) = (x-1)^{231} - (x-1)^{436}$ est donc
		\[ p(0) = (-1)^{231} - (-1)^{436} = -1 - (1) = -2, \]
	car $(-1)^n$ vaut $+1$ si $n$ est pair, et $-1$ si $n$ est impair.
}

\exe{origin=set1,difficulty=1}{
	Montrer que le polynôme $x + 1$ est combinaison linéaire des polynômes $x + 3$ et $x+5$.
	
	La combinaison trouvée est-elle unique ?
}{exe:comb-lin-poly}{
	Nous posons, pour $\alpha, \beta\in\R$ deux nombres réels,
		\[ x+1 = \alpha(x+3) + \beta(x+5) = (\alpha+\beta)x + (3\alpha + 5\beta). \]
	D'après le cours, l'ensemble $\{ 1 ; x \}$ est linéaire indépendant sur $\R$ et la combinaison linéaire de $x+1 = 1\times(x) + 1\times(1)$ est unique.
	Il suit donc que
		\[ \systeme[sort={\alpha,\beta, \gamma}]{\alpha + \beta = 1, \\ 3\alpha + 5\beta = 1.} \]
	Une façon de résoudre le système est d'écrire $\alpha = 1-\beta$ d'après la première équation, et de la substituer dans la deuxième :
		\begin{align*}
			3(1-\beta) + 5\beta &= 1, \\
			3 + 2\beta &= 1, \\
			2\beta = -2, \\
			\beta = -1.
		\end{align*}
	Par conséquent, $\alpha = 1-(-1) = 2$.
	Vérifions notre solution :
		\[ 2(x+3) + (-1)(x+5) = 2x + 6 -x-5 = x+1. \]
	La résolution du système nous ayant donné qu'une seule solution, celle-ci est unique. 
}

\exe{origin=set1,}{
	Montrer que le polynôme $x^2 - 5x + 6$ n'est pas combinaison linéaire des polynômes $x - 2$ et $x-3$.
}{exe:not-comb-lin-poly}{
	Une combinaison linéaire $\alpha(x-2) + \beta(x-3)$ avec $\alpha, \beta\in\R$ deux nombres réels ne peut pas être égale à $x^2 -5x + 6$ car le coefficient en $x^2$ de la combinaison est $0$, alors que le coefficient en $x^2$ de $x^2 -5x+6$ est 1.
}


\exe{origin=set1,difficulty=0}{
	Soit $S = \bigset{ -1 ; 0 ; 1 } \subseteq \R$ un ensemble à trois éléments.
	Considérons l'ensemble $S[x]$ des polynômes à coefficients dans $S$, ainsi que $p(x) = x^3 - x$.
	\begin{enumerate}
		\item Est-ce que $p(x) \in S[x]$ ? Justifier.
		\item Est-ce que $p$ est le polynôme nul de $S[x]$ ?
		\item Est-ce que $p$ est la fonction nulle sur $S$ ? C'est la fonction qui envoie chaque élément \mbox{de $S$ sur $0$.}
	\end{enumerate}
	\emph{Conclusion : il est donc important de distinguer polynôme et fonction polynomiale.}
}{exe:fonction-nulle-S}{
	Il s'agit de montrer que $f(x) = 0$ pour tout $x\in S$.
	Ceci se vérifie par le calcul : $f(0) = f(1) = f(-1) = 0$.
}

\exe{origin=set1,}{
	Donner la définition de l'indépendance linéaire sur $\R$ d'une famille de polynômes.
}{exe:def-LI}{
	Il existe une combinaison linéaire nulle et non triviale.
}

\exe{origin=set1,}{
	Est-ce que la famille de polynômes $\bigset{ x^2+4x - 2 ; x+2 ; 50 } \subseteq\R[x]$ est linéairement indépendante sur $\R$ ? 
}{exe:lin-indep2}{
	Oui : pour $\alpha, \beta, \gamma\in\R$ trois nombres réels, l'équation
		\[ \alpha(x^2 + 4x-2) + \beta(x+2) + \gamma(50) = 0 \]
	est équivalente à
		\[ \alpha x^2 + (4\alpha + \beta)x + (-2\alpha + 2\beta + 50\gamma) = 0. \]
	D'après l'indépendant linéaire de $\{ 1 ; x ; x^2 \}$, nous obtenons le système
		\[ \systeme[sort={\alpha,\beta, \gamma}]{\alpha = 0, \\ 4\alpha + \beta = 0, \\ -2\alpha + 2\beta + 50\gamma = 0.} \]
	La première équation donne immédiatement $\alpha =0$, qui peut disparaître des autres équations.
	La deuxième se simplifie donc en $\beta = 0$.
	La dernière se simplifie alors en $50\gamma = 0$ et donc $\gamma = 0$.
	
	Nous avons donc démontré que $\alpha=\beta=\gamma=0$ et que la seule combinaison linéaire nulle est la combinaison triviale : la famille est linéairement indépendante sur $\R$.
}

\exe{origin=set1,}{
	Est-ce que la famille de polynômes $\bigset{ x+1 ; x+2 } \subseteq\R[x]$ est libre sur $\R$ ?
}{exe:lin-indep1}{
	Oui : pour $\alpha, \beta \in \R$ deux nombres réels, l'équation
		\[ \alpha(x+1) + \beta(x+2) = 0 \]
	est équivalente à
		\[ (\alpha+\beta)x + (\alpha + 2\beta) = 0. \]
	D'après l'indépendant linéaire de $\{ 1 ; x \}$, nous obtenons le système
		\[ \systeme[sort={\alpha,\beta, \gamma}]{\alpha+\beta = 0, \\ \alpha+2\beta = 0.} \]
	En soustrayant la première équation à la deuxième, nous obtenons
		\begin{align*}
			\alpha + 2\beta - (\alpha + \beta) &= 0-0, \\
			\beta &=0.
		\end{align*}
	La première équation donne alors $\alpha = 0$ ce qui conclut.
}

\exe{origin=set1,}{
	Est-ce que la famille de polynômes $\bigset{ 2x^2+4x - 2 ; x^2 + x + 1 ; x-2 } \subseteq\R[x]$ est libre sur $\R$ ?
}{exe:lin-indep3}{
	La famille est liée car
		\[ (2x^2 + 4x - 2) - 2(x^2 + x + 1) - 2(x-2) = 0. \]
}




\exe{origin=set1,difficulty=0}{
	Démontrer que $x-1$ divise $x^8 - 1$.
}{exe:diviseur-x-m-1}{
	Comme vu dans le cours, la relation classique de somme géométrique
		\[ (1-q)(1 + q + q^2 + \dots + q^{n}) = 1-q^{n+1} \]
	permet de voir que $1-x$ divise $1-x^{n+1}$ pour tout $n\in\N$.
	
	Dans notre cas,
		\[ x^{8}-1 = (x-1) (1+x+x^2 +x^3 +x^4+x^5+x^6+x^7), \]
	que l'on vérifiera bien sûr en développant le produit !
	Nous remarquerons alors l'apparition d'une somme téléscopique : chaque terme annule le terme précédent et en ajoute un nouveau.
}

\exe{origin=set1,difficulty=1}{
	Montrer que si $f(x)$ divise $g(x)$ dans $\R[x]$, alors
	\begin{multicols}{2}
	\begin{enumerate}
		\item  $f(-x)$ divise $g(-x)$ ; et
		\item $f(x+1)$ divise $g(x+1)$.
	\end{enumerate}
	\end{multicols}
	En conclure que 
	\begin{multicols}{2}
	\begin{enumerate}
		\item  $x+1$ divise $x^8 - 1$ ; et que
		\item  $x$ divise $(x+1)^{1204} - 1$.
	\end{enumerate}
	\end{multicols}
	
	Est-ce que $x+1$ divise toujours $x^n - 1$ pour $n\geq2$ ?
}{exe:diviseur-x-p-1}{
	\begin{enumerate}
		\item
		Si $f|g$, alors $g(x) = f(x)q(x)$ avec $q\in\R[x]$.
		En remplaçant $x$ par $-x$, nous trouvons que $g(-x) = f(-x)q(-x)$, ce qui signifie que $f(-x)$ divise $g(-x)$.
		\item  
		Il s'agit de faire exactement la même chose en remplaçant $x$ par $x+1$ dans l'identité $g(x) = f(x)q(x)$.
	\end{enumerate}
	
	\begin{enumerate}
		\item
		Comme $x-1$ divise $x^{8}-1$, remplacer $x$ par $-x$ donne que $-x-1$ divise $(-x)^8 -1 = x^8 - 1$ car $8$ est pair.
		La division polynomiale ignorant les constante, $-(-x-1)$ divise aussi $x^8 -1$, comme requis.
		\item  
		Comme $x-1$ divise $x^{1204} -1$ (somme géométrique, à nouveau),remplacer $x$ par $x+1$ donne que $x$ divise $(x+1)^{1204} - 1$.
	\end{enumerate}
}

% todo : x² - 1 divise x^{120} - 1 par exemple pour les racines.

\exe{origin=set1,difficulty=1}{
	Démontrer que $x-1$ divise $x^3 - 1 + 2x - 2x^2$.
}{exe:diviseur-x-m-1-2}{
	Il s'agit de trouver $q(x)\in\R[x]$ tel que 
		\[ x^3 -1 +2x -2x^2 = (x-1)q(x). \]
	Première remarque : le degré de $q$ doit être égal à 2, et son coefficient dominant doit être égal à 1.
	Donc $q(x) = x^2 + q'(x)$ avec $\deg(q') \leq 1$.
	Ainsi
		\begin{align}
			x^3 + 1 +2x -2x^2 &= (x-1)\left (x^2 + q'(x)\right), \\
			x^3 + 1 +2x -2x^2  &= (x-1)x^2 + (x-1)q'(x), \\
			x^3 + 1 +2x -2x^2 &= x^3 -x^2 + (x-1)q'(x), \\
			1 + 2x -3x^2 &= (x-1)q'(x).
		\end{align}
	Nous avons donc réduit notre problème à la division de $1+2x-3x^2$ par $x-1$.
	Continuous en annulant une nouvelle fois le terme dominant du dividende : $\deg(q') = 1$ et son coefficient dominant doit être égale à $-3$.
	Donc $q'(x) = -3x + c$ avec $c\in\R$ et
		\begin{align*}
			1+2x-3x^2 &= (x-1)(-3x + c), \\
			1+2x-3x^2 &= -3x^2 +3x + (x-1)c, \\
			1-x &= (x-1)c.
		\end{align*}
	Bien sûr $c=-1$ fonctionne, donc finalement
		\[ q(x) = x^2 -3x - 1. \]
	Nous avons procédé à la division euclidienne de $x^3 -1+2x-2x^2$ par $x-1$ en annulant le terme dominant à chaque itération.
}

\exe{origin=set1,difficulty=2}{
	Supposons la famille $\{ p_1 ; p_2 ; p_3 \} \subseteq \R[x]$ libre.
	Montrer que les familles 
	\begin{align*}
		\mathcal{F}_1 = \bigset{ p_1 ; p_2 + p_1 ; p_3 } && \et && \mathcal{F}_2 = \bigset{ p_1 + p_2 ; p_2 + p_3 ; p_1 + p_2 + p_3 } 
	\end{align*}
	sont également libres.
}{exe:libre-to-libre}{
	Soient $\alpha, \beta, \gamma\in\R$ trois nombres réels tels que
		\[ \alpha p_1 + \beta(p_2+p_1) + \gamma p_3 = 0. \]
	Alors
		\[ (\alpha+\beta)p_1 + \beta p_2 + \gamma p_3 = 0, \]
	et, par indépendance linéaire de $\{ p_1 ; p_2 ; p_3 \}$, nous avons le système
		\[ \systeme[sort={\alpha,\beta, \gamma}]{\alpha+\beta=0, \\ \beta = 0, \\ \gamma  = 0,} \]
	dont la seule solution est bien sûr $\alpha=\beta=\gamma=0$.
	Ceci montre l'indépendance linéaire de $\mathcal{F}_1$.
	
	De façon analogue pour $\alpha, \beta, \gamma\in\R$ trois nombres réels tels que
		\[ \alpha (p_1+p_2) + \beta(p_2+p_3) + \gamma (p_1+p_2+p_3) = 0, \]
	nous avons
		\[ (\alpha + \gamma)p_1 + (\alpha + \beta + \gamma)p_2 + (\beta+\gamma)p_3 = 0. \]
	Par l'indépendance linéaire de $\{p_1 ; p_2 ; p_3\}$, nous obtenons le système
		\[ \systeme[sort={\alpha,\beta, \gamma}]{\alpha+\gamma=0, \\ \alpha+\beta+\gamma=0, \\ \beta+\gamma=0.} \]
	Le résoudre aboutit à $\alpha=\beta=\gamma=0$, ce qui montre l'indépendance linéaire de $\mathcal{F}_2$.
}


\exe{origin=set1,}{
	Vrai ou faux ? Montrer que chaque assertion est vraie ou donner un contre-exemple.
	\begin{enumerate}[itemsep=10pt]
		\item
		Si $f|g$ dans $\R[x]$, alors $\deg(f) \leq \deg(g)$.
		\item
		Si $f$ divise $g$ dans $\R[x]$, alors $(3f)$ divise aussi $g$.
		\item
		Si $f$ divise $g$ dans $\R[x]$, alors $f$ divise aussi $(3g)$.
		\item
		Soient $f, g\in\R[x]$. Alors $\deg(f+g) = \max\{\deg(f) ; \deg(g)\}$.
		\item
		Si $f|g$ dans $\R[x]$, alors $f|g^2$ dans $\R[x]$ aussi.
		\item
		Le polynôme $(1-x)^{22}$ est unitaire.
	\end{enumerate}
}{exe:TF-polynomes}{
	\begin{enumerate}
		\item
		Faux uniquement à cause du polynôme nul $g=0$. 
		Tous les polynômes divisent le polynôme nul. 
		\item
		Vrai car si $g=fq$, alors $g = (3f) (\frac13 q)$.
		\item
		Vrai car si $g=fq$, alors $3g = f (3q)$.
		\item
		Il faut et il suffit que la somme des coefficients de degré $\max\{\deg(f) ; \deg(g)\}$ soit non nulle.
		Si $g(x) = x^2 + 1$ et $f(x) = -x^2$ par exemple, alors $\deg(f+g) = 0$.
		\item
		Vrai car si $g = fq$, alors $g^2 = f(qg) = fq'$.
		\item
		Vrai : son coefficient dominant, en $x^{22}$ est donné par $(-1)^{22} = 1$ car $22$ est pair.
	\end{enumerate}
	
}


\exe{origin=set1,difficulty=1}{
	Montrer que la famille $\bigset{ 1 ; \sqrt{2} } \subseteq\R$ est linéairement indépendante sur $\Q$.
}{exe:lin-indep4}{
	Soient $p, q\in\Q$ tels que $p + \sqrt2 q = 0$.
	Supposons par contradiction que $q\neq0$.
	Alors $\sqrt2 = \frac{-p}q \in \Q$, car le produit de deux rationnels est un rationnel.
	Comme $\sqrt2$ est irrationnel, ceci est une contradiction, et $q=0$. Il suit donc que $p=0$ aussi, ce qui conclut.
}

\exe{origin=set1,difficulty=2}{
	Montrer que la famille $\bigset{ 1 ; \sqrt{2} ; \sqrt{3} } \subseteq\R$ est linéairement indépendante sur $\Q$.
}{exe:lin-indep23}{
	Supposons que $p + q\sqrt2 + r\sqrt3 = 0$ avec $p, q, r \in\Q$.
	En isolant $r\sqrt3$ et en mettant au carré, nous obtenons
		\[ p^2 + 2pq\sqrt2 + 2q^2 = 3r^2. \]
	Par indépendance linéaire de $\{ 1 ; \sqrt2 \}$ sur $\Q$, nous avons nécessairement $pq = 0$.
	Deux alternatives s'offrent donc à nous : $p=0$ ou $q=0$.
	
	Si $q=0$, alors $p+r\sqrt3 = 0$, et on montre de façon analogue à l'exercice \ref{exe:lin-indep4} que $p=r=0$ en utilisant que $\sqrt3$ est irrationnel.
	
	Si $p=0$, alors $q\sqrt2 + r\sqrt3 = 0$. 
	En multipliant par $\sqrt2$, nous obtenons $2q + r\sqrt6 = 0$.
	Le nombre $\sqrt6$ étant irrationnel lui aussi, nous obtenons $q=r=0$.
	
}

\exe{origin=set1,difficulty=3}{
	Est-ce que la famille des racines carrées de nombres premiers s'arrêtant à un certain $p$ est libre sur $\Q$ ? $$\mathcal{S}_p=\bigset{ 1 ; \sqrt{2} ; \sqrt{3} ; \sqrt{5} ; \sqrt{7} ; \sqrt{11} ; \dots ; \sqrt{p} } \subseteq \R$$
	% S for quadratic surds :)
}{exe:lin-indep-primes}{
	La réponse est oui ; la démonstration se fait par récurrence forte en généralisant la stratégie de l'exercice \ref{exe:lin-indep23}.
}


\exe{difficulty=2}{
	Soit $F_1, F_2 \subseteq \R[x]$ deux familles de polynômes.
	Nommons $V_1$ et $V_2$ l'ensemble de toutes les combinaisons linéaires de $F_1$ et $F_2$ respectivement.
	
	Montrer que, pour un polynôme $f$ combinaison linéaire de $F_1 \cup F_2$,
	\begin{align*}
		V_1 \cap V_2 = \{ 0 \} && \iff && f = v_1 + v_2 \text{ se décompose de façon unique avec $v_1\in V_1$ et $v_2\in V_2$}.
	\end{align*}
	Lorsque la décomposition est unique, nous dirons alors que $V_1$ et $V_2$ sont en \emphindex{somme directe}, notée $V_1 \oplus V_2$.
}{exe:direct-sum}{
	todo
}



\exe{difficulty=1}{
	Le but de cet exercice est de construire l'ensemble $\R[x]$ dans lequel les monômes $\{ 1 ; x ; x^2 ; \dots \}$ sont linéairement indépendants sur $\R$ à l'aide des suites.
	
	Considérons à cette fin l'ensemble $\hat{\R}^\infty$ des suites réelles ayant seulement un nombre fini de termes non nuls.
	\begin{enumerate}
		\item
		Montrer que toute suite $a\in\hat{\R}^\infty$ prend la forme
			\[ a = (a_0, a_1, \dots, a_d, 0, 0, 0,\dots). \]
	\end{enumerate}
	Nous verrons donc dans la suite un polynôme  $a_0 + a_1 x + \dots + a_d x^d$ comme une suite $a = (a_0, a_1, a_2, \dots, a_d, 0, 0, 0,\dots)$ de $\hat{\R}^\infty$.
	\begin{enumerate}[resume]
		\item
		À quelle suite correspond chaque polynôme suivant ?
		\begin{multicols}{2}
		\begin{enumerate}
			\item Le polynôme nul.
			\item Le polynôme $1$.
			\item Le polynôme $x$.
			\item Le polynôme $x^2$.
		\end{enumerate}
		\end{multicols}
		\item
		Comment définir l'addition $a+b$ ainsi que les multiplication $\alpha a$ et $a \cdot b$ pour $\alpha\in\R$ et $a, b\in\hat{\R}^\infty$ ?
		\item
		Montrer que les suites $\{ e_0 ; e_1 ; \dots ; e_d \}$ sont linéairement indépendantes sur $\R$, où
			\[ e_i = ( 0, 0, \dots, 0, \underbrace{1}_{\text{position $i+1$}}, 0, 0, \dots ). \]
		%Ces suites sont appelées \emphindex{vecteurs canoniques}.
	\end{enumerate}
}{exe:Rx-constr}{

}



%%%%%%%%%%%%%%%%%%%%%%%%
%%%%%%%%%% SET 2 %%%%%%%%%%
%%%%%%%%%%%%%%%%%%%%%%%% 



\exe{origin=set2}{
	Vrai ou faux ?
	%\begin{multicols}{2}
	\begin{enumerate}
		\item
		$1$ est une racine de $x - 1$.
		\item
		$-1$ est une racine de $x^3 - 1$. %\columnbreak
		\item
		$3$ est une racine de $(x+3)(x-2)(x+7)^4(x^2 - 1)$.
		\item
		$-1$ est une racine de $2(1-x)^{34} - (1-x)^{35}$.
	\end{enumerate}
	%\end{multicols}
}{exe:VF-is-racine}{
	\begin{enumerate}
		\item
		Vrai car $1-1 = 0$.
		\item
		Faux car $(-1)^3 - 1 = -1-1 = -2 \neq 0$.
		\item
		Faux car $(3+3)(3-2)(3+7)^4(3^2 - 1) \neq$ car aucun des facteurs n'est nul.
		\item
		Vrai car $2(1-(-1))^{34} - (1-(-1))^{35} = 2\times2^{34} - 2^{35} = 0$.
	\end{enumerate}
}


\exe{origin=set2}{
	Poser les divisions euclidiennes de
	\begin{enumerate}
		\item
		$x^3 - x$ par $x-1$ ; et de
		\item
		$3x^3 + 6x^2 - 3x + 1$ par $2x^2 + 1$ ;
	\end{enumerate}
}{exe:euclidean-division-poly}{
	\begin{enumerate}
		\item
		$\polylongdiv[style=D]{x^3 - x}{x-1}$ donc $x^3 -x = (x-1)(x^2 + x)$.
		\item
		$\polylongdiv[style=D]{3x^3 + 6x^2 - 3x + 1}{2x^2 + 1}$ donc $3x^3 + 6x^2 -3x + 1 = (2x^2 + 1)(\frac32x + 3) + (-\frac92x - 2)$.
	\end{enumerate}
}

\exe{origin=set2}{
	Donner un polynôme $p\in\R[x]$ ayant les caractéristiques suivantes :
	\begin{enumerate}
		\item
		$\deg(p) = 12$ ; 
		\item
		$1$ est une racine de $p$ de multiplicité $2$ ;
		\item
		$-4$ est une racine de $p$ de multiplicité $3$ ;
		\item
		Le coefficient constant de $p$ est nul ; et
		\item
		Le coefficient dominant de $p$ est irrationnel.
	\end{enumerate}
}{exe:poly-with-constraints}{
	Le polynôme doit être composé de $(x-1)^2$ et de $(x+4)^3$ ainsi que d'autres facteurs qui ne s'annulent pas en $1$ ni $-4$ et qui permettrons d'obtenir un degré 12.
	Comme le coefficient constant de $p$ est $p(0) = 0$, alors $0$ est également une racine de $p$ et $p$ doit admettre un facteur $x$.
	Par exemple,
		\[ f(x) = (x-1)^2(x+4)^3 x(x+1)^6. \]
	Le problème est que le coefficient dominant de $f$ est 1 (obtenu en choisissant $x$ dans chaque facteur), donc
		\[ p(x) = \sqrt2 (x-1)^2 (x+4)^3 x(x+1)^6, \]
	lui, fonctionne.
}



\exe{}{
	Considérons l'équation quadratique $3x^2-7x+4= 0$ pour $x\in\R$.
	Montrer que $1$ est racine de l'équation et en déduire l'autre racine à l'aide d'une division euclidienne.
}{exe:euclide-poly}{
	En remplaçant $x$ par 1, on trouve bien $3-7+4 = 0$.
	La division euclidienne de $3x^2 -7x + 4$ par $x-1$ donne
	
	$\polylongdiv[style=D]{3x^2 -7x + 4}{x-1}$ et donc $3x^2 - 7x + 4 = (x-1)(3x-4)$.
	
	Il suit donc que $3x-4 = 0 \iff x= \frac43$ est l'autre racine de l'équation.
}


\exe{}{
	Trouver une racine triviale (évidente) puis factoriser le polynôme $f(x) = 7x^2 - 10x + 3$ pour trouver sa deuxième racine.
}{exe:euclide-poly2}{
	La somme des coefficients est nulle, donc $f(1) = 0$.
	La division euclidienne de $7x^2 - 10x + 3$ par $x-1$ donne
	
	$\polylongdiv[style=D]{7x^2 -10x + 3}{x-1}$ et donc $7x^2 - 10x + 3 = (x-1)(7x -3)$.
	
	L'autre racine de $f$ est donc $7x - 3 = 0 \iff x = \frac37$.
}


\exe{}{
	Considérons le polynôme $f(x) = 21x^3 - 7x^2 - 9x + 3$.
	Montrer par le calcul que $\frac13$ est une racine de $f$.
	En déduire les autres racines de $f$.
	
	\emph{Indication : poser la division euclidienne de $f$ par $3x-1$ plutôt que par $x-\frac13$ pour simplifier les calculs.}
}{exe:euclide-poly3}{
	\[ \polylongdiv[style=D]{21x^3 -7x^2 -9x + 3}{3x-1} \]

	Par division euclidienne, $f(x) = (3x-1)(7x^2-3)$, donc les racines sont $\frac13$ et $\pm\sqrt{\frac37}$.
}


\exe{}{
	Montrer que $x^2 - 1$ divise $x^{2026}-1$.
}{exe:xsq-m-1}{
	Le nombre 1 est racine de $x^{2026} - 1$, donc $x-1$ divise $x^{2026}-1$ dans $\R[x]$ :
		\[ x^{2026} - 1 = (x-1)q(x). \]
	De plus, $-1$ est racine de $x^{2026}-1$, donc $(-1-1)q(-1) = 0 \iff q(-1) = 0$.
	Par suite, $-1$ est racine de $q(x)$ et $x+1$ divise $q(x)$ :
		\[ x^{2026} - 1 = (x-1)(x+1)\tilde{q}(x) = (x^2 -1)\tilde{q}(x). \]
}


% NOT USED
\exe{}{
	Considérons le polynôme $f(x) = 4x^3 - 6x^2 - 2x + 3$.
	Montrer, par le calcul, que $\frac32$ est une racine de $f$.
	En déduire les autres racines de $f$.
}{exe:8}{
		\[ (2x - 3)(2x^2 - 1) = 4x^3 - 2x -6x^2 + 3. \]
	On utilise que le produit est nul si et seulement si un des deux facteurs est nul.
	Or l'équation $2x^2 - 1 = 0$ est équivalente à $x^2 = \dfrac12$.
	En prenant la racine et en utilisant que $\sqrt{x^2} = |x|$, on obtient
		\[ |x| = \sqrt{\dfrac12} = \dfrac1{\sqrt2}. \]
	Les deux valeurs $\pm \dfrac1{\sqrt2}$ (notation $\pm$ pour \og plus ou moins \fg) sont donc solutions.
	On en déduit que l'ensemble des solutions de \eqref{eq:2} est $\left\{ \dfrac32 ; \pm \dfrac1{\sqrt2} \right\} = \left\{ \dfrac32 ; \dfrac1{\sqrt2} ; - \dfrac1{\sqrt2} \right\}$

}

\exe{}{
	Soient $k, n\in\N$.
	Montrer que la $k$-ième dérivée de $x^n$ est égale à $0$ si $n<k$ et $\frac{n!}{(n-k)!} x^{n-k}$ sinon.
	En particulier, $(x^n)^{(n)} = n!$.
}{exe:deriv-xn}{
	Si $n\geq k$, alors
		\[ (x^n)^{(k)} = n \times (x^{n-1})^{(k-1)} = n\times(n-1)\times \dots \times(n-k+1)x^{n-k} = \frac{n!}{(n-k)!}x^{n-k}. \]
	Si $n < k$, alors, d'après ce qui a été vu au-dessus,
		\[ (x^n)^{(n)} = n! \]
	est un polynôme constant. Dériver une nouvelle fois le rend nul, donc
		\[ (x^n)^{(k)} = 0. \]
}


\exe{difficulty=1}{
	Montrer le corollaire du théorème de Taylor vu en classe : 
	soit $f\in\R[x]$ un polynôme à coefficients réels de degré $d$ et $\alpha\in\R$ un réel quelconque.
	Alors
		\[ f(x) = \sum_{k=0}^{d} \frac{f^{(k)}(\alpha)}{k!} (x-\alpha)^k.\]
	\emph{Indication : appliquer le théorème de Taylor au polynôme $f(x+\alpha)$.}
}{exe:taylor-poly-2}{
	Il suffit d'appliquer le théorème de Taylor au polynôme $g(x) = f(x+a)$.
	Comme $g'(x) = f'(x+\alpha)$, et $g^{(k)}(x) = f^{(k)}(x+\alpha)$ pour tout $k\in\N^*$,
	nous avons
		\[ f(x+\alpha) = g(x) = \sum_{k=0}^d \frac{g^{(k)}(0)}{k!} x^k = \sum_{k=0}^d \frac{f^{(k)}(\alpha)}{k!} x^k. \]
	En remplaçant $x$ par $x-\alpha$, nous obtenons bien
		\[ f(x) = \sum_{k=0}^d \frac{f^{(k)}(\alpha)}{k!} (x-\alpha)^k. \]
}



\exe{}{
	Considérons le polynôme
		\[ f(x) = (x-3)^2 - (x-3)^4 + 3(x-3)^5. \]
	Calculer $f(3), f'(3), f''(3), f'''(3), f^{(4)}(3)$, et $f^{(5)}(3)$ à l'aide du théorème de Taylor (exercice \ref{exe:taylor-poly-2}).
}{exe:f-shift-derivatives}{
	D'après le théorème de Taylor,
		\[ f(x) = \sum_{k=0}^{d} \frac{f^{(k)}(3)}{k!} (x-3)^k = 0 + 0(x-3) + (x-3)^2 + 0(x-3)^3 - (x-3)^4 + \frac1{12}(x-3)^5. \]
	Par identification des coefficients, nous avons
	\begin{align*}
		\frac{f(3)}{0!} = 0, && \frac{f'(3)}{1!} = 0, && \frac{f''(3)}{2!} = 1, && \frac{f'''(3)}{3!} = 0, && \frac{f^{(4)}(3)}{4!} = -1, && \frac{f^{(5)}(3)}{5!} = \frac1{12}.
	\end{align*}
	et donc
	\begin{align*}
		f(3) = 0, && f'(3) = 0, && f''(3) = 2, && f'''(3) = 0, && f^{(4)}(3) = -24, && \frac{f^{(5)}(3)}{5!} = 10.
	\end{align*}
}


\exe{difficulty=2}{
	Soit $f\in\R[x]$ un polynôme à coefficients réels et $\alpha\in\R$ un réel quelconque.
	%
	Montrer, à l'aide du théorème de Taylor, que $\alpha$ est une racine de multiplicité $m$ si et seulement si 
		\begin{align*}
		 f(\alpha) = f'(\alpha) = \cdots = f^{(m-1)}(\alpha) = 0 && \et && f^{(m)}(\alpha) \neq 0.
		\end{align*}
}{exe:mult-from-diff}{
	\underline{Direction $\impliedby$.}
	
	Supposons que
	\begin{align*}
		 f(\alpha) = f'(\alpha) = \cdots = f^{(m-1)}(\alpha) = 0 && \et && f^{(m)}(\alpha) \neq 0.
	\end{align*}
	Déjà, $d = \deg(f) > m$ car la $m$-ième dérivée de $f$ n'est pas nulle.
	
	Alors, d'après le théorème de Taylor,
	\begin{align*}
		f(x) &= \sum_{k=0}^{d} \frac{f^{(k)}(\alpha)}{k!} (x-\alpha)^k = \sum_{k=m}^{d} \frac{f^{(k)}(\alpha)}{k!} (x-\alpha)^k 
		\\
		&= (x-\alpha)^m  \sum_{k=m}^{d} \frac{f^{(k)}(\alpha)}{k!} (x-\alpha)^{k-m}
		\\
		&= (x-\alpha)^m q(x),
	\end{align*}
	donc $(x-\alpha)^m$ divise bien $f(x)$.
	Il faut désormais montrer que $(x-\alpha)^{m+1}$ ne divise pas $f(x)$ pour conclure.
	Or,
	\begin{align*}
		f(x) &= x-\alpha)^m  \sum_{k=m}^{d} \frac{f^{(k)}(\alpha)}{k!} (x-\alpha)^{k-m}
		\\
		&= \frac{f^{(m)}(\alpha)}{m!} (x-\alpha)^m + (x-\alpha)^{m+1} \sum_{k=m+1}^{d} \frac{f^{(k)}(\alpha)}{k!} (x-\alpha)^{k-m-1}
		\\
		&= (x-\alpha)^{m+1} \tilde{q}(x) + \frac{f^{(m)}(\alpha)}{m!} (x-\alpha)^m 
	\end{align*}
	est exactement la division euclidienne de $f(x)$ par $(x-\alpha)^{m+1}$ car $\deg(x-\alpha)^m = m < \deg(x-\alpha)^{m+1} = m+1$.
	
	Or, le reste n'est pas nul car $f^{(m)}(\alpha) \neq 0$ par hypothèse, donc nous concluons que $(x-\alpha)^{m+1}$ ne divise pas $f(x)$, comme requis.
	
	\underline{Direction $\implies$.}
	
	Supposons que $\alpha\in\R$ soit une racine de $f$ de multiplicité $m\in\N^*$.
	Considérons l'unique entier naturel $k\in\N^*$ vérifiant
	\begin{align*}
		 f(\alpha) = f'(\alpha) = \cdots = f^{(k-1)}(\alpha) = 0 && \et && f^{(k)}(\alpha) \neq 0.
	\end{align*}
	Il existe car $f(\alpha) = 0$ parce que $\alpha$ est une racine de $f$, et il est unique parce qu'il correspond au premier moment où la suite des $f^{(i)}(\alpha)$ est non nulle.
	
	D'après la direction $\impliedby$ déjà montrée, $\alpha$ est une racine de $f$ multiplicité $k$.
	Par conséquent, $k=m$, et nous avons bien
	\begin{align*}
		 f(\alpha) = f'(\alpha) = \cdots = f^{(m-1)}(\alpha) = 0 && \et && f^{(m)}(\alpha) \neq 0.
	\end{align*}
	
%%% useless lol
%
%	Si $\alpha\in\R$ est une racine de $f$ de multiplicité $m$, alors $(x-\alpha)^m$ divise $f$ mais $(x-\alpha)^{m+1}$ ne divise pas $f$.
%	Le théorème de Taylor donne la division euclidienne de $f(x)$ par $(x-\alpha)^m$ dont le reste doit donc être nul :
%	\begin{align*}
%		f(x) &= \sum_{k=0}^{d} \frac{f^{(k)}(\alpha)}{k!} (x-\alpha)^k 
%		\\
%		&= \sum_{k=0}^{m-1} \frac{f^{(k)}(\alpha)}{k!} (x-\alpha)^k + (x-\alpha)^m q(x)
%		&= (x-\alpha)^m q(x) + 0
%	\end{align*}
%	Par conséquent, la somme $\sum_{k=0}^{m-1} \frac{f^{(k)}(\alpha)}{k!} (x-\alpha)^k $ est égale au polynôme nul.
%	
%	Or, le degré de $(x-\alpha)^k$ est égal à $k$, donc seul $(x-\alpha)^{m-1}$ contribue en facteur $x^{m-1}$ dont le coefficient doit être nul.
%	Il suit donc que $ \frac{f^{(m-1)}(\alpha)}{(m-1!} = 0 \iff f^{(m-1)}(\alpha) = 0$, et que la somme s'arrête au terme $k=m-2$ :
%		\[ \sum_{k=0}^{m-2} \frac{f^{(k)}(\alpha)}{k!} (x-\alpha)^k  = 0. \]
%	De façon analogue, seul $(x-\alpha)^{m-2}$ contribue au coefficient en $x^{m-2}$, et $f^{(m-2)}(\alpha) = 0$.
%	
%	En répétant le raisonnement, nous obtenons bien que $f(\alpha) = f'(\alpha) = \cdots = f^{(m-1)}(\alpha) = 0$.
%	Si $f^{(m)}(\alpha)$ était nul, nous aurions que $\alpha$ est une racine de multiplicité au moins $m+1$ d'après la direction $\impliedby$ déjà traitée.
%	Ce n'est pas le cas donc $f^{(m)}(\alpha)= 0$.
}

\exe{difficulty=2}{
	Considérons le \emphindex{symbole de Pochhammer} défini par
		\begin{align*}
			(x)_0 = 1, && \et && (x)_n = x(x-1)(x-2)\cdots(x-n+1) \text{ pour $n\in\N^*$, }
		\end{align*}
	ainsi que l'opérateur de différence $\Delta$ qui associe à chaque $f\in\R[x]$ le polynôme 
		\[ (\Delta f)(x) = f(x+1)-f(x). \]
	
	Montrer que, pour tout $n\in\N^*$, l'identité suivante est vraie dans $\R[x]$.
		\[ \Delta(x)_n = n\cdot(x)_{n-1} \]
	\emph{Indication : quelles sont les racines de $(x)_{n-1}$ et de $\Delta(x)_n$ ?}
}{exe:Pochhammer}{
	Remarquons que les racines de $(x)_n$ sont exactement $0 ; 1 ; 2 ; \dots ; n-1$ pour $n\geq1$.
	De plus, si $x^*$ et $x^* +1$ sont deux racines de $f$, alors $x^*$ est racine de $\Delta(f)$ car $\Delta(f)(x^*) = f(x^* + 1)- f(x^*) = 0-0=0$.

	Pour $n\geq 2$, posons $f(x) = \Delta(x)_n$ et identifions toutes les racines de $f$.
	D'après la remarque, les racines $\{ 0; 1 ; \dots ; n-2 \}$ sont des racines de $f$.
	Il suit donc que $x(x-1)\dots(x-n+2) = (x)_{n-1}$ divise $f(x)$ dans $\R[x]$.
	
	Regardons les termes dominants de $f(x)$, en $x^n$ et $x^{n-1}$ :
	\begin{align*}
		f(x) &= (x+1)(x+2)\cdots(x+n) - x(x+1)\cdots(x+n-1) \\
		&= x^n + (1 + 2 + 3 + \dots + n)x^{n-1} + \underbrace{...}_{\text{puissances de $x$ inférieures}} \\
		&- \left[ x^n + (0 + 1 + 2 + \dots + n-1)x^{n-1} + \underbrace{...}_{\text{puissances de $x$ inférieures}} \right] \\
		&= 0x^n + nx^{n-1} + \underbrace{...}_{\text{puissances de $x$ inférieures}}.
	\end{align*}
	Par conséquent, $\deg(f) = n-1$, mais $(x)_{n-1}$ divise $f(x)$ et est aussi de degré $n-1$.
	Il suit que $f(x) = c \cdot (x)_{n-1}$ pour une constante $c\in\R$.
	Comme $(x)_{n-1}$ est unitaire, mais que $f(x) = nx^{n-1} + ...$, nous obtenons $c=n$.
	
}

\exe{difficulty=0}{
	Considérons la famille $F$ suivante.
		\[ F = \bigset{ (x-1)(x-2)(x-3) ; x(x-2)(x-3) ; x(x-1)(x-3) ; x(x-1)(x-2) } \]
	\begin{enumerate}
		\item
		Lister les racines de chaque polynôme de $F$.
		\item
		\underline{Sans développer} les expressions, montrer que la famille $F$ est linéairement \mbox{indépendante sur $\R$.}
	\end{enumerate}
}{exe:unique-roots}{
	\begin{enumerate}
		\item
		\begin{enumerate}[label=Polynôme \arabic* :]
			\item
			Les racines sont $1 ; 2; 3$.
			\item
			Les racines sont $0 ; 2; 3$.
			\item
			Les racines sont $0 ; 1 ; 3$.
			\item
			Les racines sont $0 ; 1 ; 2$.
		\end{enumerate}
		\item
		Soient $\alpha, \beta, \gamma, \delta \in\R$ quatre nombres réels tels que
			\[ \alpha(x-1)(x-2)(x-3) + \beta x(x-2)(x-3) + \gamma x(x-1)(x-3) + \delta x(x-1)(x-2) = 0. \]
		En évaluant en $x=0$, nous obtenons
			\[ \alpha \times (-6) = 0 \iff \alpha = 0. \]
		En évaluant en $x=1$, nous obtenons
			\[ \beta \times 2 = 0 \iff \beta = 2. \]
		En évaluant en $x=2$, nous obtenons
			\[ \gamma \times (-2) = 0 \iff \gamma = 0. \]
		En évaluant en $x=3$, nous obtenons finalement
			\[ \delta \times 6 = 0 \iff \delta = 0. \]
		La seule combinaison linéaire nulle est la combinaison linéaire triviale : les polynômes sont linéairement indépendants sur $\R$.
	\end{enumerate}
}

\exe{difficulty=1}{
	Le but de cet exercice est de généraliser l'exercice \ref{exe:unique-roots}.
	
	Considérons une famille de polynômes $\bigset{ f_1, f_2, \dots, f_n} \subseteq \R[x]$ telle que pour tout $i=1, 2, \dots n$, il existe un nombre $a_i\in\R$ qui est racine de $f_j$ pour tout $j\neq i$, mais n'est pas racine de $f_i$.
	Montrer que la famille est linéairement indépendante sur $\R$.
}{exe:unique-roots2}{
	Soient $\alpha_1, \dots, \alpha_n \in\R$ des nombres réels tels que
		\[ \alpha_1 f_1 + \alpha_2 f_2 + \dots + \alpha_n f_n = 0. \]
	En évaluant en $x=a_i$, nous obtenons
		\[ \alpha_i f_i(a_i) = 0 \iff \alpha_i = 0, \]
	car $f_j(a_i) = 0$ pour tout $j\neq i$, et $f_i(a_i) \neq 0$.
	Ceci étant vrai pour tout $i=1 ; 2 ; \dots ; n$, nous concluons que la seule combinaison linéaire nulle est la combinaison triviale : la famille est libre sur $\R$.
}


\exe{origin=set2,}{
	Vrai ou faux ? Montrer que chaque assertion est vraie ou donner un contre-exemple.
	\begin{enumerate}[itemsep=10pt]
		%\item
		%Soient $f, g\in\R[x]$. Alors $\deg(fg) = \deg(f)+\deg(g)$.
		\item
		Un polynôme constant de $\R[x]$ est de degré nul.
		\item
		$124(x+1)$ divise $2x^5 -x^4 - 3x^3 + x^2 -3x - 4$ dans $\R[x]$.
		\item
		Si $f$ divise $g$ dans $\R[x]$, alors toutes les racines de $f$ sont des racines de $g$.
		\item
		Si toutes les racines de $f$ sont des racines de $g$, alors $f$ divise $g$ dans $\R[x]$.
		\item
		Un polynôme non nul de $\R[x]$ de degré $n\in\N$ admet au plus $n$ racines distinctes.
	\end{enumerate}
}{exe:TF-polynomes2}{
	\begin{enumerate}
		%\item
		%Vrai car le coefficient dominant de $fg$ est le produit des coefficients dominants de $f$ et de $g$.
		%Le cas $f=0$ est vrai par convention que $\deg(0) = \minfty$.
		\item
		Faux à cause du polynôme nul de degré $\minfty$ par convention.
		\item
		Vrai car $-1$ est racine du deuxième polynôme.
		\item
		Vrai car si $g = fq$, et que $a$ est une racine de $f$, alors $g(a) = f(a)q(a) = 0q(a) = 0$.
		\item
		Faux car la multiplicité d'une racine de $f$ peut être trop élevée.
		Un contre-exemple est donc $f(x) = (x-1)^2$ et $g(x) = x-1$ : toutes les racines de $f$ (donc que $1$) sont des racines de $g$, mais $f$ ne divise pas $g$ car son degré est strictement supérieur et $g$ est non nul.
		\item
		Vrai. Par contradiction, si $\alpha_1, \dots, \alpha_{n+1}$ sont des racines distinctes de $f$, alors $\prod_{i=1}^{n+1} (x-\alpha_i)$ divise $f$ dans $\R[x]$.
		Or, le produit est de degré $n+1$ et $f$ est de degré $n$, ce qui n'est pas possible car $f$ est supposé non nul.
	\end{enumerate}
	
}

\exe{}{
	Montrer qu'une fonction polynomiale de $\R[x]$ de degré 3 admet au moins une racine réelle.
}{exe:deg3-root}{
	C'est une application du théorème des valeurs intermédiaires.
}

\exe{difficulty=2}{
	Montrer la \emphindex{règle du produit} : pour deux polynômes $f, g\in\R[x]$,
		\[ (fg)' = f'g + g'f. \]
}{exe:deriv-prod}{
	À gauche, le coefficient en $x^n$ de $fg$ est égal à
		\[ \sum_{i+j=n} f_i g_j. \]
	Le coefficient en $x^{n-1}$ de $fg$ est donc égal à
		\[ n\sum_{i+j=n} f_i g_j. \]
	À droite, le coefficient en $x^{n-1}$ de $f'g$ est égal à
		\[ \sum_{i+j=n-1} f'_ig_j = \sum_{i+j=n-1} (i+1)f_{i+1}g_j = \sum_{i+j=n} if_ig_j. \]
	Le coefficient en $x^{n-1}$ de $g'f$ est égal à
		\[ \sum_{i+j=n-1} f_ig'_j = \sum_{i+j=n-1} f_i (j+1)g_{j+1} = \sum_{i+j=n} f_ijg_j. \]
	Le coefficient en $x^{n-1}$ de $f'g + g'f$ est donc égal à
		\[ \sum_{i+j=n} f_ijg_j +  \sum_{i+j=n} if_ig_j = \sum_{i+j=n} (i+j)f_ig_j = n\sum_{i+j=n} f_i g_j. \]
	Les coefficients de $(fg)'$ et $f'g+g'f$ sont donc égaux et les polynômes sont égaux.
}


\exe{difficulty=1}{
	Soit $f$ une fonction dérivable en 0.
	Montrer que 
		\[ f(x) = f(0) + f'(0)x + g(x), \]
	où $g(x) = o(x)$ dans le sens suivant : 
		\[ \lim_{x \to 0} \frac{g(x)}x = 0. \]
}{exe:taylor-C1}{
	Posons $g(x) = f(x) - f(0) - f'(0)x$.
	Pour $x\neq0$, nous avons
		\[ \frac{g(x)}x = \frac{f(x) - f(0)}{x-0} - f'(0). \]
	Lorsque $x$ tend vers $0$,  $\frac{f(x) - f(0)}{x-0}$ tend vers $f'(0)$ par définition de $f'(0)$.
	Il suit donc que $\frac{g(x)}x$ tend vers 0 comme requis.
	
	Remarque : si $f\in C^2(\R)$, il s'avère que $g(x) = \O(x^2)$, c'est-à-dire que $\frac{g(x)}{x^2}$ reste borné lorsque $x$ tend vers $0$ (ce qui est plus fort que $g=o(x)$).
	Nous aurons en fait 
		\[ f(x) = f(0) + f'(0)x + \frac{f''(0)}2 x^2 + o(x^2). \]
}

\exe{difficulty=2}{
	Soit $f\in\R[x]$ un polynôme à coefficients réels.
	Supposons que $z\in\C$ soit une racine complexe non réelle de $f$.
	\begin{enumerate}
		\item
		Montrer que le conjugué $\overline{z}$ est aussi un racine de $f$.
		\item
		En déduire que $x^2 - 2\RE(z) x + |z|^2$ divise $f$ dans $\R[x]$.
	\end{enumerate}
	Un polynôme est dit \emphindex{irréductible} s'il est non constant et si ses seuls diviseurs sont les constantes et les multiples de lui-même.
	\begin{enumerate}[resume]
		\item
		En admettant le théorème fondamental de l'algèbre, montrer que les polynômes irréductibles de $\R[x]$ sont nécessairement de degré 1 ou 2.
	\end{enumerate}
}{exe:irr-in-R}{
	\begin{enumerate}
		\item
		Comme $f$ est à coefficients réels et que $\overline{z}^n = \overline{z^n}$ pour tout $n\in\N$ et $z\in\C$, nous avons que $0 = \overline{0} = \overline{f(z)} = f(\overline{z})$.
		\item
		Comme la racine $z$ n'est pas réelle, $\overline{z} \neq z$ et nous avons trouvé deux racines de $f$ dans $\C$.
		Le polynôme $g(x) = (x-z)(x-\overline{z})$ divise donc $f$ dans $\C[x]$.
		
		Or $g$ est à coefficients réels, car $g(x) = x^2 - 2\RE(z) x + |z|^2$.
		La division euclidienne de $f$ par $g$ dans $\R[x]$ a donc un sens et nous allons montrer qu'elle correspond à celle dans $\C[x]$.
		Soient $q, r\in\R[x]$ tels que
			\[ f = gq + r, \]
		avec $\deg(r) < \deg(g) = 2$.
		Alors $r$ est une fonction affine réelle qui s'annule en $z$ et en $\overline{z} \neq z$, car $r= f-gq$.
		Or, dans $\R$ comme dans $\C$, l'équation $ax+b = 0$ n'a plus de deux solutions que dans le cas $a=b=0$ (si $a\neq 0$, alors la solution est unique, et si $a=0, b\neq0$, alors il n'y a pas de solution).
		Il suit donc que $r = 0$ et que $g$ divise $f$ dans $\R[x]$.
		\item
		Un polynôme $f\in\R[x]$ de degré supérieur ou égal à 3 admettra nécessairement une racine complexe d'après le théorème fondamental de l'algèbre.
		Si cette racine $r$ est réelle, alors $x-r$ est un diviseur propre de $f(x)$ qui est réductible.
		Si cette racine n'est pas réelle, alors $(x-r)(x-\overline{r})$ est un diviseur propre de $f(x)$ qui est également réductible.
	\end{enumerate}
}

\exe{}{
	Calculer la division euclidienne de $a(x) = 2x^3 + 3x^2 - 5$ par $b(x) = x+11$ et en déduire la valeur de $a(-11)$.
	
	% c'e'st pas ouf car c'est plus compliqué de calculer la division que l'image
	 % car on trouve (x+1)(2x^2 -19x + 209).
	 % peut--être avec de meilleurs choix ?
	 % c'est tiré de Ullman, Aho, Hopcroft.
}{exe:div-for-image}{
	$\polylongdiv[style=D]{2x^3 + 3x^2 - 5}{x+11}$
}


\exe{difficulty=1}{
	%Soit $f\in C^1(\R)$ une fonction continûment dérivable.
	Soit $f:\R\rightarrow\R$ une fonction dérivable en 0.
	Montrer que 
		\[ f(x) = f(0) + f'(0)x + g(x), \]
	où $g(x) = o(x)$ dans le sens suivant : 
		\[ \lim_{x \to 0} \frac{g(x)}x = 0. \]
}{exe:Taylor-C1}{
	Posons $g(x) = f(x) - f(0) - f'(0)x$.
	Pour $x\neq0$, nous avons
		\[ \frac{g(x)}x = \frac{f(x) - f(0)}{x-0} - f'(0). \]
	Lorsque $x$ tend vers $0$,  $\frac{f(x) - f(0)}{x-0}$ tend vers $f'(0)$ par définition de $f'(0)$.
	Il suit donc que $\frac{g(x)}x$ tend vers 0 comme requis.
	
	Remarque : si $f\in C^2(\R)$, il s'avère que $g(x) = \O(x^2)$, c'est-à-dire que $\frac{g(x)}{x^2}$ reste borné lorsque $x$ tend vers $0$ (ce qui est plus fort que $g=o(x)$).
	Nous aurons en fait 
		\[ f(x) = f(0) + f'(0)x + \frac{f''(0)}2 x^2 + o(x^2). \]
}


%%%%%%%%%%%%%%%%%%%%%%%%
%%%%%%%%%% SET 3 %%%%%%%%%%
%%%%%%%%%%%%%%%%%%%%%%%% 



\exe{}{
	Calculer les combinaisons linéaires dans $\R^n$ suivantes.
	\begin{align*}
		u = \begin{pmatrix} 7 \\ -2 \\ 4 \end{pmatrix}  + \begin{pmatrix} -2 \\ -1 \\ 4 \end{pmatrix}
		&&
		v = -2\begin{pmatrix}2 \\ -\pi \\ \frac12 \end{pmatrix}
		&&
		w = 3\begin{pmatrix}2\\-4\\3\end{pmatrix} + \begin{pmatrix}2\\4\\-1\end{pmatrix}
		&&
		z = 2\begin{pmatrix} 1\\-1\\3\\-7\end{pmatrix} - 10\begin{pmatrix}-1\\-1\\2\\3\end{pmatrix}
	\end{align*}
}{exe:calcul-vectoriel}{
	\begin{multicols}{4}
	\begin{enumerate}
		\item
		$\begin{pmatrix} 7 \\ -2 \\ 4 \end{pmatrix}  + \begin{pmatrix} -2 \\ -1 \\ 4 \end{pmatrix} = \begin{pmatrix} 9 \\ -1 \\ 0 \end{pmatrix}$
		\item
		$-2\begin{pmatrix}2 \\ -\pi \\ \frac12 \end{pmatrix} = \begin{pmatrix}-4 \\ 2\pi \\ -1 \end{pmatrix} \in \R^3$
		\item
		$3\begin{pmatrix}2\\-4\\3\end{pmatrix} + \begin{pmatrix}2\\4\\-1\end{pmatrix} = \begin{pmatrix}3\times2+2\\3\times(-4)+4\\3\times3+(-1)\end{pmatrix} =\begin{pmatrix}8\\-8\\8\end{pmatrix}$
		\item
		$2\begin{pmatrix} 1\\-1\\3\\-7\end{pmatrix} - 10\begin{pmatrix}-1\\-1\\2\\3\end{pmatrix} = \begin{pmatrix}12\\8\\-14\\-44\end{pmatrix}$
	\end{enumerate}
	\end{multicols}
}


\exe{}{
	Pour chaque paire de vecteurs suivante, sont-ils orthogonaux ?
	\begin{multicols}{3}
	\begin{enumerate}
		\item $\pvec12$ et $\pvec{-1}{-2}$
		\item $\pvec12$ et $\pvec{-2}{1}$
		\item $\pvec{-6}{-3}$ et $\pvec{1}{-2}$
		\item $\pvec{-6}{-3}$ et $\pvec{-1}{2}$
		\item $\begin{pmatrix}2 \\ 4 \\ -3 \end{pmatrix}$ et $\begin{pmatrix}5 \\ -1 \\ -2 \end{pmatrix}$
		\item $\begin{pmatrix}2 \\ 4 \\ -3 \end{pmatrix}$ et $\begin{pmatrix}-5 \\ 1 \\ 2 \end{pmatrix}$
	\end{enumerate}
	\end{multicols}
}{exe:is-orth}{
	todo
}



\exe{}{
	Est-ce que la famille de vecteurs suivante est linéairement indépendante sur $\R$ ?
		\[ \left\{ \begin{pmatrix} 2 \\ -1 \\ 3 \end{pmatrix} ; \begin{pmatrix} 0 \\ 2 \\ 4 \end{pmatrix} ; \begin{pmatrix} 1 \\ 0 \\ 1 \end{pmatrix} \right\} \]
}{exe:lin-indep-vec}{

}

\exe{difficulty=1}{
	Soit $a\in\R$. Donner un vecteur orthogonal à $\pvec1a$.
}{exe:vec-orth}{
	todo
}

\exe{difficulty=0}{
	Si $u, v\in\R^n$ sont orthogonaux, montrer que $u$ et $-v$ le sont aussi puis que $u$ et $\lambda v$ le sont aussi pour tout $\lambda\in\R$.
	%Interpréter géométriquement.
}{exe:vec-orth-lin}{
	todo
}




\exe{}{
	Calculer la distance entre les points $P(3;4;2;-8)$ et $Q(-2;2;-3;2)$ dans $\R^4$.
}{exe:dist-4d}{
	todo
}

\exe{}{
	Soient $u, v$ deux vecteurs de $\R^n$. Montrer l'identité remarquable
		\[ \norm{u+v}^2 = \norm{u}^2 + \norm{v}^2 + 2\scal{u, v}. \]
}{exe:id-rem-scal}{
	todo
}

\exe{difficulty=1}{
	Soient $v_1, \dots, v_n \in \R^n$ une famille orthonormale : $\japb{v_i, v_j} = 0$ si $i\neq j$, et $1$ sinon.
	\begin{enumerate}
		\item 
		Montrer l'identité de Pythagore 
			\[ \norm{v_1 + v_2 + \dots + v_n}^2 = \norm{v_1}^2 + \norm{v_2}^2 + \dots + \norm{v_n}^2 \]
		\item
		Montrer que la famille est linéairement indépendante sur $\R$.
	\end{enumerate}
}{exe:base-orthonormale}{

}




\exe{, difficulty=0}{
	Pour chaque système d'équations d'inconnues $x, y \in \R$, donner un système équivalent tel que les coefficients multipliant $x$ soient opposés (c'est-à-dire leur somme soit nulle).
	\setlength{\columnsep}{1cm}
	\begin{multicols}{2}
	\begin{enumerate}[label=\roman*), itemsep=20pt]
		\item $\systeme{x + 3y = 3, \\  -2x + 2y = 2.}$ % \iff \systeme{\,,\,}$
		\item $\systeme{x - 7y = 1, \\  -3x + 5y = 13.}$ % \iff \systeme{\,,\,}$
		\item $\systeme{2x - y = 3, \\  -x + 2y = -5.}$ % \iff \systeme{\,,\,}$
		\item $\systeme{4x  - y = -4, \\   x + 5y = 2.}$ % \iff  \systeme{\,,\,}$
		\item $\systeme{-2x - 2y = -5, \\  3x + 7y = 10.}$ % \iff\systeme{\,,\,}$
		\item $\systeme{7x - y = 10, \\  5x + 2y = 1.}$ % \iff\systeme{\,,\,}$
	\end{enumerate}
	\end{multicols}
}{exe:systeme2}{
	\setlength{\columnsep}{1cm}
	\begin{multicols}{2}
	\begin{enumerate}[label=\roman*), itemsep=20pt]
		\item $\systeme{x + 3y = 3, \\  -2x + 2y = 2.}$ % \iff \systeme{2x+6y=6, \\ -2x+2y=2.}$
		\item $\systeme{x - 7y = 1, \\  -3x + 5y = 13.}$ % \iff \systeme{3x - 21y = 3, \\  -3x + 5y = 13.}$
		\item $\systeme{2x - y = 3, \\  -x + 2y = -5.}$ % \iff \systeme{2x - y = 3, \\  -2x + 4y = -10.}$
		\item $\systeme{4x  - y = -4, \\   x + 5y = 2.}$ % \iff  \systeme{4x  - y = -4, \\   -4x - 20y = -8.}$
		\item $\systeme{-2x - 2y = -5, \\  3x + 7y = 10.}$ % \iff\systeme{-6x - 6y = -15, \\  6x + 14y = 20.}$
		\item $\systeme{7x - y = 10, \\  5x + 2y = 1.}$ % \iff \systeme{35x - 5y = 50, \\  -35x - 14y = -7.}$
	\end{enumerate}
	\end{multicols}
}

\exe{}{
	Pour chaque système de l'exercice \ref{exe:2}, 
		\begin{enumerate}
			\item combiner les équations pour trouver $y$ ; 
			\item trouver $x$ ; 
			\item vérifier que le couple $(x ;y)$ obtenu soit bien solution du système initial.
		\end{enumerate}
}{exe:systeme3}{ \, \\
	\begin{enumerate}[label=Système \roman*) :, itemsep=20pt, leftmargin=80pt]
		\item La somme des deux équations donne
			\[ 8y = 8 \iff y = 1. \]
		En remplaçant $y$ par $1$ dans la première équation, on trouve $x + 3 = 3$, et donc $x =0$.
			\[ (x ; y ) = (0; 1). \]
		
		\item La somme des deux équations donne
			\[ -16y = 16 \iff y = -1. \]
		En remplaçant $y$ par $-1$ dans la première équation, on trouve $x + 7 = 1$, et donc $x =-6$.
			\[ (x ; y ) = (-6; -1). \]
		
		\item La somme des deux équations donne
			\[ 3y = -7 \iff y = -\frac73. \]
		En remplaçant $y$ par $-\frac73$ dans la première équation, on trouve $2x +\frac73 = 3$, et donc $x =\frac13$.
			\[ (x ; y ) = \left(\frac13; -\frac73 \right). \]
		
		\item La somme des deux équations donne
			\[ -21y = -12 \iff y = \frac{12}{21} = \frac47. \]
		En remplaçant $y$ par $ \frac47$ dans la deuxième équation, on trouve $x +\frac{20}7 = 2$, et donc $x =-\frac67$.
			\[ (x ; y ) = \left(-\frac67; \frac47 \right). \]
		
		\item La somme des deux équations donne
			\[ 8y = 5 \iff y = \frac58. \]
		En remplaçant $y$ par $\frac58$ dans la deuxième équation, on trouve $3x +\frac{35}8 = 10 \iff x =\frac{45}{24} = \frac{15}8$.
			\[ (x ; y ) = \left(\frac{15}8; \frac58 \right). \]
		
		\item La somme des deux équations donne
			\[ -19y = 43 \iff y = -\frac{43}{19}. \]
		En remplaçant $y$ par $-\frac{43}{19}$ dans la première équation, on trouve $7x +\frac{43}{19}= 10 \iff x =\frac{147}{133} = \frac{21}{19}$.
			\[ (x ; y ) = \left( \frac{21}{19};  -\frac{43}{19} \right). \]
	\end{enumerate}
}

\exe{}{
	Donner les solutions réelles $(x; y)$ des systèmes suivantes.
	
	\begin{multicols}{2}
	\begin{enumerate}[label=\roman*), itemsep=20pt]
		\item $\systeme{2x+4y=0, \\ 2x -2y = 2.}$
		\item $\systeme{-6x+3y=-3, \\ x -3y = -2.}$
		\item $\systeme{x+y = 2, \\ x+y=2.}$
		\item $\systeme{8x-4y = 6+2x-y, \\ 7x+12y=-5+y.}$
		\item $\systeme{x-y = 1, \\ -x+y=10.}$
		\item $\systeme{\frac12 x - \frac23 y= -1, \\ \frac15 x + \frac72 y= 5.}$
		\item $\systeme{2x - 8y = 2, \\ -4x+16y=-1.}$
		\item $\systeme{5x + y = -2x-1-y, \\ 8x = -1+2y.}$
		\item $\systeme{2y + 12x = -3,\\ -6x=y + \frac32.}$
	\end{enumerate}
	\end{multicols}
}{exe:systeme4}{\, \\
	
	\begin{enumerate}[label=\roman*), itemsep=20pt]
		\item $(x ; y) = (-2 ; 1)$
		\item $(x ; y) = (1;1)$
		\item Les deux équations sont redondantes : on a en fait qu'une seule contrainte pour deux variables et donc un degré liberté. Pour chaque choix de $x$, on peut trouver un $y$ tel que $(x;y)$ soit solution.
		L'ensemble des solutions est
			\[ \{ (x ; y) \text{ tq. } x+y=2, x, y \in \R \} = \{ (x ; y) \text{ tq. } y = -x + 2, \text{ où $x$ parcourt $\R$} \} = \Ccal_f, \]
		où $f(x) = -x+2$.
		\item $(x ; y) = \left(\frac{17}{29} ; -\frac{24}{29}\right)$
		\item Les deux équations sont contradictoires et aucune solution n'existe.
		\item $(x ; y) = \left(-\frac{10}{113} ; \frac{162}{113}\right)$
		\item Les deux équations sont contradictoires et aucune solution n'existe.
		\item $(x ; y) = \left(-\frac{2}{15} ; -\frac{1}{30}\right)$
		\item Les deux équations sont redondantes : on a en fait qu'une seule contrainte pour deux variables et donc un degré liberté. Pour chaque choix de $x$, on peut trouver un $y$ tel que $(x;y)$ soit solution.
		L'ensemble des solutions est
			\[ \left\{ (x ; y) \text{ tq. } -6x -y = \frac32, x, y \in \R \right\} = \left\{ (x ; y) \text{ tq. } y= -6x - \frac32 \text{ où $x$ parcourt $\R$} \right\} = \Ccal_f, \]
		où $f(x) =-6x - \frac32$.
	\end{enumerate}
}



\exe{}{
	Soient $u = \begin{pmatrix}1\\3\\-4\end{pmatrix}$ et $v = \begin{pmatrix} 1 \\ -1 \\ 0\end{pmatrix}$.
	Quel $\lambda\in\R$ choisir pour que les vecteurs $v$ et $u-\lambda v$ soient orthogonaux ?
	
	En déduire le point de plus proche de $P(1;3;-4)$ appartenant à la droite $\R v$.
}{exe:Fourier-coeff}{

}

\exe{}{
	Trouver le point le plus proche de $P(2;-1;0)$ appartenant à $(2; 1; -1) + \R \begin{pmatrix}2\\0\\-1\end{pmatrix}$.
}{exe:proj-orth-R3}{

}

\exe{}{
	Soient $u = \begin{pmatrix}34\\17\\-17\end{pmatrix}$, $v = \begin{pmatrix} 0 \\ -1 \\ -1\end{pmatrix}$ et  $w = \begin{pmatrix} 2 \\ -2 \\ 1\end{pmatrix}$.
	Quels $\lambda, \mu\in\R$ choisir pour que fois les vecteurs $v$ et $u-\lambda v-\mu w$ soient orthogonaux et les vecteurs $w$ et $u-\lambda v -\mu w$ soient orthogonaux ?
	Poser un système de deux équations à deux inconnues puis le résoudre.
	
	En déduire le point de plus proche de $P(34;17;-17)$ appartenant au plan $\R v + \R w$.
}{exe:proj-orth-plan}{

}

\exe{difficulty=1}{
	Montrer que la famille $\mathcal{F}$ suivante est orthonormée.
		\[ \mathcal{F} = \left\{ \frac1{\sqrt6}\begin{pmatrix} 2 \\ 2 \\ 1 \end{pmatrix} ; \frac1{\sqrt6}\begin{pmatrix} -1 \\ 2 \\ -2 \end{pmatrix} ; \frac1{\sqrt6}\begin{pmatrix} -2 \\ 1 \\ 2 \end{pmatrix} \right\} \]
	Exprimer le point
		\[ P(5 ;  5 ; -2 ) \]
	comme combinaison linéaire des vecteurs de $\mathcal{F}$.
	
	\emph{Indication : ignorer les $\frac1{\sqrt6}$ pour simplifier les calculs et les rajouter à la fin.}
		
	En déduire $Q$, le point le plus proche de $P$ appartenant au plan $\R\begin{pmatrix} 2 \\ 2 \\ 1 \end{pmatrix} + \R\begin{pmatrix} -2 \\ 1 \\ 2 \end{pmatrix}$.
}{exe:proj-orthonormal}{

}

\exe{difficulty=2}{
	Soient $b_1^*, b_2, b_3 \in\R^3$ trois vecteurs linéairement indépendants.
	\begin{enumerate}
		\item
		Montrer que 
			\[ b_2^* = b_2 - \frac{\japb{b_2, b_1^*}}{\japb{b_1^*, b_1^*}} b_1^* \]
		est orthogonal à $b_1^*$.
		\item
		Montrer que 
			\[ b_3^* = b_3 - \frac{\japb{b_3, b_1^*}}{\japb{b_1^*, b_1^*}} b_1^* - \frac{\japb{b_3, b_2^*}}{\japb{b_2^*, b_2^*}} b_2^* \]
		est orthogonal à $b_1^*$ et $b_2^*$.
	\end{enumerate}
	\emph{Remarque : cet algorithme peut continuer et porte le nom de \emphindex{procédé de Gram-Schmidt}.}
}{exe:Gram-Schmidt}{
	todo
}



\exe{difficulty=1}{
	Soient $u, v$ deux vecteurs non nuls de $\R^2$ et $\theta$ l'angle dirigé de $u$ vers $v$.
	%En adaptant la preuve du théorème \ref{thm:ab-orth}, montrer que
	Montrer que
		\[ \scal{u, v} = \norm{u}\norm{v}\cos(\theta). \]
}{exe:prod-scal-2d}{
	todo
}


\exe{difficulty=1}{
	Soient $u, v\in\R^n$ deux vecteurs linéairement indépendants.
	\begin{enumerate}
		\item
		Développer la fonction quadratique $f(t) = \norm{u + tv}^2$.
		\item
		Que dire du discriminant du polynôme quadratique $f$ ?
		\item
		En déduire l'\emphindex{inégalité de Cauchy-Schwarz} :
			\[ \abs{\japb{u, v}} \leq \norm{u} \norm{v}. \]
		\item
		Montrer l'égalité suivante pour n'importe quels réels $x_1, \dots, x_n \in \R$.
			\[ |x_1| + \dots + |x_n| \leq \sqrt{n} \sqrt{x_1^2 + \dots + x_n^2} \]
	\end{enumerate}
}{exe:Cauchy-Schwarz}{
	\begin{enumerate}
		\item
		\begin{align*}
			f(t) &= \norm{u + tv}^2, \\
			&= \norm{u}^2 + \norm{tv}^2 + 2\japb{u, tv}, \\
			&= t^2 \norm{v}^2 + 2t \japb{u, v} + \norm{u}^2.
		\end{align*}
		\item
		Le polynôme est toujours positif ou nul car c'est le carré d'une norme.
		Son discriminant est donc négatif ou nul car $f$ n'admet qu'une ou aucune racine.
		\item
		\begin{align*}
			\Delta = b^2 - 4ac &\leq 0, \\
			\left(2\japb{u, v}\right)^2 - 4\norm{v}^2\norm{u}^2 &\leq 0, \\
			4\left(\japb{u, v}^2 - \norm{v}^2\norm{u}^2\right) &\leq 0, \\
			\japb{u, v}^2 &\leq  \norm{v}^2\norm{u}^2, \\
			\abs{\japb{u, v}} &\leq \norm{u} \norm{v}.
		\end{align*}
		\item
		En prenant $u$ le vecteur dont toutes les coordonnées sont 1 et $v$ dont les coordonnées sont $|x_1|, \dots, |x_n|$, l'inégalité de Cauchy-Schwarz donne le résultat attendu.
	\end{enumerate}
}

\exe{difficulty=2}{
	Considérons une série numérique $(a_n)_{n\in\N}$ telle que $\sum_{n\in\N} a_n^2 < \infty$.
	Montrer, à l'aide de l'inégalité de Cauchy-Schwarz (exercice \ref{exe:Cauchy-Schwarz}) que
		\[ \sum_{n\in\N} \frac{|a_n|}{n^s} < \infty \]
	pour tout $s>\frac12$.
}{exe:Cauchy-Schwarz-application}{

}









